\documentclass[12pt,letterpaper]{article}
\usepackage{graphicx,textcomp}
\usepackage{natbib}
\usepackage{setspace}
\usepackage{fullpage}
\usepackage{color}
\usepackage[reqno]{amsmath}
\usepackage{amsthm}
\usepackage{fancyvrb}
\usepackage{amssymb,enumerate}
\usepackage[all]{xy}
\usepackage{endnotes}
\usepackage{lscape}
\newtheorem{com}{Comment}
\usepackage{float}
\usepackage{hyperref}
\newtheorem{lem} {Lemma}
\newtheorem{prop}{Proposition}
\newtheorem{thm}{Theorem}
\newtheorem{defn}{Definition}
\newtheorem{cor}{Corollary}
\newtheorem{obs}{Observation}
\usepackage[compact]{titlesec}
\usepackage{dcolumn}
\usepackage{tikz}
\usetikzlibrary{arrows}
\usepackage{multirow}
\usepackage{xcolor}
\newcolumntype{.}{D{.}{.}{-1}}
\newcolumntype{d}[1]{D{.}{.}{#1}}
\definecolor{light-gray}{gray}{0.65}
\usepackage{url}
\usepackage{listings}
\usepackage{color}

\definecolor{codegreen}{rgb}{0,0.6,0}
\definecolor{codegray}{rgb}{0.5,0.5,0.5}
\definecolor{codepurple}{rgb}{0.58,0,0.82}
\definecolor{backcolour}{rgb}{0.95,0.95,0.92}

\lstdefinestyle{mystyle}{
	backgroundcolor=\color{backcolour},   
	commentstyle=\color{codegreen},
	keywordstyle=\color{magenta},
	numberstyle=\tiny\color{codegray},
	stringstyle=\color{codepurple},
	basicstyle=\footnotesize,
	breakatwhitespace=false,         
	breaklines=true,                 
	captionpos=b,                    
	keepspaces=true,                 
	numbers=left,                    
	numbersep=5pt,                  
	showspaces=false,                
	showstringspaces=false,
	showtabs=false,                  
	tabsize=2
}
\lstset{style=mystyle}
\newcommand{\Sref}[1]{Section~\ref{#1}}
\newtheorem{hyp}{Hypothesis}

\title{Problem Set 3}
\date{Due: March 25, 2026}
\author{Applied Stats II}


\begin{document}
	\maketitle
	\section*{Instructions}
	\begin{itemize}
	\item Please show your work! You may lose points by simply writing in the answer. If the problem requires you to execute commands in \texttt{R}, please include the code you used to get your answers. Please also include the \texttt{.R} file that contains your code. If you are not sure if work needs to be shown for a particular problem, please ask.
\item Your homework should be submitted electronically on GitHub in \texttt{.pdf} form.
\item This problem set is due before 23:59 on Wednesday March 25, 2026. No late assignments will be accepted.
	\end{itemize}

	\vspace{.25cm}
\section*{Question 1}
\vspace{.25cm}
\noindent We are interested in how governments' management of public resources impacts economic prosperity. Our data come from \href{https://www.researchgate.net/profile/Adam_Przeworski/publication/240357392_Classifying_Political_Regimes/links/0deec532194849aefa000000/Classifying-Political-Regimes.pdf}{Alvarez, Cheibub, Limongi, and Przeworski (1996)} and is labelled \texttt{gdpChange.csv} on GitHub. The dataset covers 135 countries observed between 1950 or the year of independence or the first year forwhich data on economic growth are available ("entry year"), and 1990 or the last year for which data on economic growth are available ("exit year"). The unit of analysis is a particular country during a particular year, for a total $>$ 3,500 observations. 

\begin{itemize}
	\item
	Response variable: 
	\begin{itemize}
		\item \texttt{GDPWdiff}: Difference in GDP between year $t$ and $t-1$. Possible categories include: "positive", "negative", or "no change"
	\end{itemize}
	\item
	Explanatory variables: 
	\begin{itemize}
		\item
		\texttt{REG}: 1=Democracy; 0=Non-Democracy
		\item
		\texttt{OIL}: 1=if the average ratio of fuel exports to total exports in 1984-86 exceeded 50\%; 0= otherwise
	\end{itemize}
	
\end{itemize}
\newpage
\noindent Please answer the following questions:
\begin{enumerate}
\item Construct and interpret an unordered multinomial logit with \texttt{GDPWdiff} as the output and "no change" as the reference category, including the estimated cutoff points and coefficients.
		\lstinputlisting[language=R, firstline=49, lastline=56]{PS3.R} 
	\begin{table}[!htbp] \centering 
		\caption{Unordered Multinominal Logit} 
		\label{} 
		\begin{tabular}{@{\extracolsep{5pt}}lcc} 
			\\[-1.8ex]\hline 
			\hline \\[-1.8ex] 
			& \multicolumn{2}{c}{\textit{Dependent variable:}} \\ 
			\cline{2-3} 
			\\[-1.8ex] & negative & positive \\ 
			\\[-1.8ex] & (1) & (2)\\ 
			\hline \\[-1.8ex] 
			REG & 1.379$^{*}$ & 1.769$^{**}$ \\ 
			& (0.769) & (0.767) \\ 
			& & \\ 
			OIL & 4.784 & 4.576 \\ 
			& (6.885) & (6.885) \\ 
			& & \\ 
			Constant & 3.805$^{***}$ & 4.534$^{***}$ \\ 
			& (0.271) & (0.269) \\ 
			& & \\ 
			\hline \\[-1.8ex] 
			Akaike Inf. Crit. & 4,690.770 & 4,690.770 \\ 
			\hline 
			\hline \\[-1.8ex] 
			\textit{Note:}  & \multicolumn{2}{r}{$^{*}$p$<$0.1; $^{**}$p$<$0.05; $^{***}$p$<$0.01} \\ 
		\end{tabular} 
	\end{table} 
\begin{itemize}
		\item In a given country, there is an increase in the baseline odds (no change in GDP) that the difference in GDP will be negative by 1.379 times when a country is a democracy. 
		\item  In a given country, there is an increase in the baseline odds (no change in GDP) that the difference in GDP will be positive by 1.769 times when a country is a democracy. 
		\item  In a given country, there is an increase in the baseline odds (no change in GDP) that the difference in GDP will be negative by 4.784 times given the average ratio of fuel exports in 1984-86 exceeded 50 percent. 
		\item  In a given country, there is an increase in the baseline odds (no change in GDP) that the difference in GDP will be postive by 4.576 times given the average ratio of fuel exports in 1984-86 exceeded 50 percent. 
	\end{itemize}
	\item Construct and interpret an ordered multinomial logit with \texttt{GDPWdiff} as the outcome variable, including the estimated cutoff points and coefficients.
\lstinputlisting[language=R, firstline=65, lastline=69]{PS3.R} 
	\begin{table}[!htbp] \centering 
		\caption{Ordered Multinominal Logit} 
		\label{} 
		\begin{tabular}{@{\extracolsep{5pt}}lc} 
			\\[-1.8ex]\hline 
			\hline \\[-1.8ex] 
			& \multicolumn{1}{c}{\textit{Dependent variable:}} \\ 
			\cline{2-2} 
			\\[-1.8ex] & GDPWdiff\_cat \\ 
			\hline \\[-1.8ex] 
			REG & 0.410$^{***}$ \\ 
			& (0.075) \\ 
			& \\ 
			OIL & $-$0.179 \\ 
			& (0.115) \\ 
			& \\ 
			\hline \\[-1.8ex] 
			Observations & 3,721 \\ 
			\hline 
			\hline \\[-1.8ex] 
			\textit{Note:}  & \multicolumn{1}{r}{$^{*}$p$<$0.1; $^{**}$p$<$0.05; $^{***}$p$<$0.01} \\ 
		\end{tabular} 
	\end{table} 
	\begin{table}[!htbp] \centering 
		\caption{Odds Ratio and CI}
		\label{} 
		\begin{tabular}{@{\extracolsep{5pt}} cccc} 
			\\[-1.8ex]\hline 
			\hline \\[-1.8ex] 
			& OR & 2.5 \% & 97.5 \% \\ 
			\hline \\[-1.8ex] 
			REG & $1.507$ & $1.301$ & $1.747$ \\ 
			OIL & $0.836$ & $0.668$ & $1.051$ \\ 
			\hline \\[-1.8ex] 
		\end{tabular} 
	\end{table} 
	\begin{itemize}
		\item For a country that is a democracy, the log odds of being in a higher GDP change category are 1.507 times higher compared to a country that isn't a democracy, holding all other variables constant. 
		\item IF the coefficent was statisically significant: For a country that had the average ratio of fuel exports to total exports in 1984-86 that exceeded 50 percent, the log odd of being in a higher GDP change category are 0.836 times higher compared to a country that did not, holding all other variables constant. 
	\end{itemize}
\end{enumerate}

\newpage
\section*{Question 2} 
\vspace{.25cm}

\noindent Consider the data set \texttt{MexicoMuniData.csv}, which includes municipal-level information from Mexico. The outcome of interest is the number of times the winning PAN presidential candidate in 2006 (\texttt{PAN.visits.06}) visited a district leading up to the 2009 federal elections, which is a count. Our main predictor of interest is whether the district was highly contested, or whether it was not (the PAN or their opponents have electoral security) in the previous federal elections during 2000 (\texttt{competitive.district}), which is binary (1=close/swing district, 0="safe seat"). We also include \texttt{marginality.06} (a measure of poverty) and \texttt{PAN.governor.06} (a dummy for whether the state has a PAN-affiliated governor) as additional control variables. 

\begin{enumerate}
	\item [(a)]
	Run a Poisson regression because the outcome is a count variable. Is there evidence that PAN presidential candidates visit swing districts more? Provide a test statistic and p-value.
		\begin{itemize}
		\item Null hypothesis: There is no difference in PAN visits in whether the district is a close/swing or "safe seat". 
		\item Alternative hypothesis: There are more PAN visits to a swing district
	\end{itemize}	
	\lstinputlisting[language=R, firstline=84, lastline=90]{PS3.R} 
\begin{table}[!htbp] \centering 
	\small
	\caption{} 
	\label{} 
	\begin{tabular}{@{\extracolsep{5pt}}lc} 
		\\[-1.8ex]\hline 
		\hline \\[-1.8ex] 
		& \multicolumn{1}{c}{\textit{Dependent variable:}} \\ 
		\cline{2-2} 
		\\[-1.8ex] & PAN.visits.06 \\ 
		\hline \\[-1.8ex] 
		competitive.district & $-$0.081 \\ 
		& (0.171) \\ 
		& \\ 
		marginality.06 & $-$2.080$^{***}$ \\ 
		& (0.117) \\ 
		& \\ 
		PAN.governor.06 & $-$0.312$^{*}$ \\ 
		& (0.167) \\ 
		& \\ 
		Constant & $-$3.810$^{***}$ \\ 
		& (0.222) \\ 
		& \\ 
		\hline \\[-1.8ex] 
		Observations & 2,407 \\ 
		Log Likelihood & $-$645.606 \\ 
		Akaike Inf. Crit. & 1,299.213 \\ 
		\hline 
		\hline \\[-1.8ex] 
		\textit{Note:}  & \multicolumn{1}{r}{$^{*}$p$<$0.1; $^{**}$p$<$0.05; $^{***}$p$<$0.01} \\ 
	\end{tabular} 
\end{table} 
\newpage
\begin{description}
	\item The p value of 0.63 for the coefficent competitive.district is greater than 0.05, therefore, we fail to reject the null hypothesis that there is no difference between PAN visits in whether the district is a close/swing or "safe seat". There is insufficent evidence to support the alternative hypothesis that there are more PAN visits in swing districts. 
\end{description}
	\item [(b)]
	Interpret the \texttt{marginality.06} and \texttt{PAN.governor.06} coefficients.
	\begin{itemize}
		\item An increase of marginality by 1 unit  decreases the number of visits by PAN presidential canidate by a multiplicative factor of $\exp -2.080$ $\approx$ 0.125. 
		\item When a district has a PAN-affiliated governor, it decreases the number of visits by PAN presidential canidate by a multiplicative factor of $\exp -0.312$ $\approx$ 0.73.
	\end{itemize}
	\item [(c)]
	Provide the estimated mean number of visits from the winning PAN presidential candidate for a hypothetical district that was competitive (\texttt{competitive.district}=1), had an average poverty level (\texttt{marginality.06} = 0), and a PAN governor (\texttt{PAN.governor.06}=1).
\begin{equation}
	log(\lambda) = -3.810 - 0.081(competitive.district) - 2.080(marginality.06) - 0.312(PAN.governor.06)
\end{equation}
\begin{equation}
	log(\lambda) = -3.810 - 0.081(1) - 2.080(0) - 0.312(1)
\end{equation}
\begin{equation}
	log(\lambda) = - 4.203
\end{equation}
\begin{description}
	\item The mean number of visits from the winning PAN presidential canidate for a district that was competitive, had an average poverty level and a PAN governor is -4.203. Since the problem is dealing with visits, it is impossible to visit a negative amount of times to a district so it can be said that the number of visits is 0. 
\end{description}	
\end{enumerate}

\end{document}
